\documentclass[11pt]{article}
\usepackage[a4paper,margin=27mm]{geometry}
\usepackage[T1]{fontenc}
\usepackage{lmodern,microtype,amsmath,amssymb,amsthm,mathtools,enumitem,booktabs}
\usepackage[hidelinks,hypertexnames=false]{hyperref}
\newtheorem{theorem}{Theorem}
\newtheorem{lemma}[theorem]{Lemma}
\newtheorem{proposition}[theorem]{Proposition}
\newtheorem{corollary}[theorem]{Corollary}
\theoremstyle{remark}\newtheorem{remark}[theorem]{Remark}
\newcommand{\R}{\mathbb R}
\newcommand{\Dscr}{\mathcal D}
\newcommand{\Escr}{\mathcal E}
\newcommand{\supp}{\operatorname{supp}}
\newcommand{\one}{\mathbf 1}
\title{Beyond Additive Guidance Rankings\\[3pt]\large What incomparabilities and non-additive comparisons can achieve}
\author{}\date{9 October 2026}
\begin{document}\maketitle

\noindent
This investigation gives three results. First, the obstruction in v3 extends
to any ranking whose comparisons can be respected by an additive guidance
loss. This rules out introducing incomparabilities by intersecting additive
rankings, even infinitely many. Second, genuinely non-additive independent
rankings do exist: an explicit construction satisfies every requirement
except (S), including full coverage and the condition on ties. Third, two
natural ways of going further---comparisons required to hold on every finite
sample of problems, and absolute priority for avoiding bad advice---cannot
satisfy (S).

\emph{Full coverage with all the requirements, including (S), remains open.}
The positive constructions below are partial results, not solutions to that
remaining problem. The only theorem imported from v3 is stated explicitly
in Section~\ref{sec:setup}; all new arguments are supplied here.

\section{Setting and the result from v3}\label{sec:setup}

There are $n\ge2$ worlds and $\Delta=\{p\ge0:\sum_i p_i=1\}$.
Models are measurable acceptance sets $\Dscr\subseteq\R^n$, identified modulo
Lebesgue-null sets. We integrate using Gaussian measure $\mu$ and count each
binary problem once, on $H=\{g:g_1>0\}$.
A model gives one of four recommendations at $\{g,-g\}$: accept $g$ only,
accept neither, accept $-g$ only, or accept both. At world $i$, the correct
recommendation is the one with positive payoff there. Coordinate hyperplanes
are null and can be ignored. The guidance ranks are
\[
 r_i(g,\Dscr)=0,1,2,3
 \quad\text{for good, silent, bad, contradictory advice, respectively.}
\]
The background comparison is pointwise improvement of these ranks, a.e.;
it is strict if some improvement occurs on a positive-measure set.
The ideal model is $D_i=\{g:g_i>0\}$.

A coherent model excludes zero, contains every nonzero nonnegative gamble,
and is closed under addition and positive scaling. Modulo null sets it is
\[
 \Dscr_M=\{g:\min_{m\in M}m\cdot g>0\}
       \cup\{g\ge0:g\ne0\}
\]
for a nonempty compact convex $M\subseteq\Delta$.
The vacuous model is $\Dscr_\Delta$.

At world $i$, $\Dscr\preceq_i\Escr$ means that $\Dscr$ is at least as good.
Dominance requires weak improvement everywhere and strict improvement
somewhere. Strict dominance requires strict improvement everywhere.
A model is competitive on $S\ne\varnothing$ if, against every model, it is
at least as good at some world in $S$.

All rankings here evaluate every measurable model, giving (E). We require:
\begin{itemize}[leftmargin=*,itemsep=3pt]
\item \textbf{(G):} preserve weak and strict guidance improvements.
\item \textbf{(IB):} changing shared advice outside any measurable collection
of binary problems cannot change a comparison. Equivalently, for $A=-A$,
replacing the common part on $A^c$ by any other common part preserves both
directions of every comparison.
\item \textbf{(LA):} at least one coherent model is undominated.
\item \textbf{(A):} each coherent model is undominated under some ranking
in the class under consideration.
\item \textbf{(C):} a model tied at every world with a coherent undominated
model is itself coherent and undominated, modulo null sets.
\item \textbf{(S):} there is a selection $p\mapsto C_p$ of coherent
undominated models, competitive on $\supp p$, such that
$\{p:C_p\preceq_i\Dscr\}$ is closed for every $i,\Dscr$.
\end{itemize}

An \emph{additive guidance loss} means a finite integral
\begin{equation}
 L_i(\Dscr)=\int_H\ell_i(g,a_\Dscr(g))\,d\mu,
 \qquad 0=\ell_i(\mathrm{good})<\ell_i(0)
 <\ell_i(\mathrm{bad})<\ell_i(\times),
 \label{eq:additive}
\end{equation}
with an integrable envelope. The penalties may depend on the problem and
the world. A $p$-Bayes model minimizes $\sum_i p_iL_i$.
We use the following result, proved in Appendix C of
\emph{Full Coverage by Homothetic Credal Families v3}:
\begin{quote}
If $n\ge3$ and an additive guidance loss has a coherent $p$-Bayes model
at every prior, no $\Dscr_M$ with three simplex vertices in $M$ is Bayes
at any prior. Moreover, every model not strictly dominated everywhere
by an additive loss is Bayes at some prior.
\end{quote}
The second assertion uses convexity of the closure of attainable loss
vectors, obtained by splitting sets of binary problems, and a separating
hyperplane. Neither assertion assumes uniqueness of Bayes models.

\section{Why removing additive comparisons cannot suffice}

The next lemma explains why (S) makes incomparabilities difficult to use.
It requires enough comparisons to put a coherent model below every model.
An undominated model must consequently share the comparisons of a selected
coherent model.

\begin{lemma}\label{lem:complete-class}
For any preorders satisfying (LA) and (S):
\begin{enumerate}[label=(\alph*)]
\item Every model is weakly dominated at all worlds by some selected $C_p$.
\item Every undominated model is tied at every world with some $C_p$,
and is therefore competitive on $\supp p$.
\end{enumerate}
\end{lemma}
\begin{proof}
For a fixed $\Dscr$, put $A_i=\{p:C_p\preceq_i\Dscr\}$.
These sets are closed, and each $p$ lies in an $A_i$ with $p_i>0$.
Their intersection is nonempty. For completeness, let $d_i(p)$ be the
distance from $p$ to $A_i$ and apply Brouwer's theorem to
\[
 F_i(p)=\frac{p_i+d_i(p)}{1+\sum_jd_j(p)}.
\]
At a fixed point, $p_i\sum_jd_j(p)=d_i(p)$. One index has both $p_i>0$
and $d_i(p)=0$, so every distance is zero. A point of the intersection
gives (a). If $\Dscr$ is undominated, every comparison in (a) must be a tie.
Transitivity then transfers competitiveness from $C_p$ to $\Dscr$.
\end{proof}

Say that $L_i$ \emph{weakly respects} a ranking if
\begin{equation}
 \Dscr\preceq_i\Escr\quad\Longrightarrow\quad L_i(\Dscr)\le L_i(\Escr).
 \label{eq:extension}
\end{equation}
Strict ordinal comparisons are allowed to become numerical ties. No
converse implication is assumed.

\begin{theorem}[Obstruction to additive weak extensions]\label{thm:extension}
Let $n\ge3$. Suppose a ranking on all models satisfies (LA) and (S), and
its comparisons are weakly respected by additive guidance losses at every
world. Then every coherent target whose credal set contains three simplex
vertices is dominated by a coherent model under the original ranking.
In particular, the vacuous model is not undominated.
\end{theorem}
\begin{proof}
Every $p$-Bayes model for $L$ has a coherent weak dominator under the original
ranking, by Lemma~\ref{lem:complete-class}(a). By \eqref{eq:extension}, this
dominator is also $p$-Bayes for $L$. Thus $L$ has coherent Bayes models at
every prior.

If a stated target were undominated under the original ranking,
Lemma~\ref{lem:complete-class}(b) would make it competitive there, and hence
competitive for $L$. It could not be strictly dominated everywhere for $L$,
so would be $p$-Bayes for some prior. This contradicts the result from v3.
The target is therefore dominated. Applying
Lemma~\ref{lem:complete-class}(a) to a dominator gives a coherent dominator.
\end{proof}

The conclusion is ordinary dominance in the original ranking; we do not
claim strict improvement at every world there. The theorem does not need
(IB), or completeness. In particular, deleting any number of comparisons
from an additive guidance ranking cannot solve the coverage problem while
retaining (S). Neither can a lexicographic refinement whose first criterion
is already an additive \emph{strict} guidance loss.

\begin{corollary}[Even arbitrarily many additive criteria do not help]
Suppose
\[
 \Dscr\preceq_i\Escr\quad\Longleftrightarrow\quad
 L_i^\alpha(\Dscr)\le L_i^\alpha(\Escr)\quad\text{for every }\alpha,
\]
where the $L_i^\alpha$ are finite integrals with locally weakly ordered
penalties: good $\le$ silence $\le$ bad $\le$ contradiction.
If the combined ranking satisfies strict (G), it has an additive guidance
weak extension. Thus it cannot satisfy (LA), (S), and leave the vacuous
model undominated when $n\ge3$.
\end{corollary}
\begin{proof}
Fix a world and one adjacent pair of guidance ranks. The nonnegative local
penalty differences, as $\alpha$ varies, must collectively be positive on
almost every problem. Otherwise improving that adjacent pair on a common
positive-measure set would tie under every criterion, violating strict (G).

A family of measurable sets in a finite measure space has a countable
subfamily with the same essential union. To see this, take finite unions
whose measures tend to the supremum of the measures of all finite unions,
and unite the selected sets. Any further set can add only a null set.
Apply this to the sets on which the local differences are positive, for
each of the three adjacent rank pairs and each world.

We obtain countably many criteria whose positive weighted sum has strictly
ordered local penalties a.e. Subtract each criterion's good-advice penalty
and choose positive coefficients small enough that the sum of their
integrable envelopes is integrable, for example $2^{-k}/(1+\int h_kd\mu)$.
The resulting loss has form \eqref{eq:additive}. Every comparison of the
intersection is respected term by term, hence by this sum.
\end{proof}

\section{A genuinely non-additive partial ranking}

The preceding theorem does not settle the ordinal problem. Here is an
explicit ranking with (G), (IB), and no additive guidance weak extension.
It even satisfies (E), (LA), full (A), and (C). Its failure is (S).

Start with the finite order $Q$ on subsets of $\{1,\ldots,6\}$ in
Appendix~\ref{app:finite}. It strictly respects inclusion and satisfies
\begin{equation}
 A\le_Q B\quad\Longleftrightarrow\quad A\setminus B\le_Q B\setminus A.
 \label{eq:finite-cancellation}
\end{equation}
Yet it admits no strictly positive additive weak extension. Four of its
comparisons suffice to prove this:
\begin{equation}
 136<_Q2,\qquad26<_Q134,\qquad23<_Q56,\qquad145<_Q236.
 \label{eq:four-witness}
\end{equation}
Here $136$ abbreviates $\{1,3,6\}$. Summing the four differences
``right indicator minus left indicator'' gives $-e_1$.
Additive weights respecting these comparisons would therefore satisfy
$0\le-w_1$, contradicting $w_1>0$.

For a vector $x\in\{0,1,2,3\}^6$, form the three sets
$T_s(x)=\{j:x_j\ge s\}$, $s=3,2,1$.
Order vectors by comparing $T_3$, then $T_2$, then $T_1$, using $Q$ at each
stage. Write $\rho(x)$ for an integer rank representing this finite order.
It is strictly increasing in the penalty
ranks, and changing coordinates shared by two vectors leaves their
comparison unchanged, by \eqref{eq:finite-cancellation}.

Partition $H$ into six radial regions $B_1,\ldots,B_6$ of positive measure,
for example $j-1<\|g\|<j$ for $j\le5$ and $\|g\|>5$ for $j=6$.
Each region contains a positive-measure set of unanimously positive gambles.
Draw one problem from each region, independently, using normalized Gaussian
measure. For a tuple $\boldsymbol g=(g^1,\ldots,g^6)$ put
\[
 \mathcal L_i(\Dscr)(\boldsymbol g)
 =\rho\bigl(r_i(g^1,\Dscr),\ldots,r_i(g^6,\Dscr)\bigr).
\]
The penalty is a function of six problems, and rankings compare these
functions almost everywhere:
\begin{equation}
 \Dscr\preceq_i\Escr\quad\Longleftrightarrow\quad
 \mathcal L_i(\Dscr)(\boldsymbol g)
 \le\mathcal L_i(\Escr)(\boldsymbol g)
 \quad\text{for almost every tuple}.
 \label{eq:tuple-ranking}
\end{equation}

\begin{proposition}\label{prop:genuine}
The ranking \eqref{eq:tuple-ranking} satisfies (E), (LA), (A), (C), (G),
and (IB). Every coherent model is undominated under this single ranking.
At no world does it have an additive guidance weak extension.
\end{proposition}
\begin{proof}
Reflexivity and transitivity hold pointwise. Coordinatewise monotonicity of
$\rho$ gives (G), including strictness: an improvement on a positive-measure
part of any region affects a positive-measure set of tuples.
At each tuple, changing shared advice changes only coordinates shared by
the two rank vectors. The finite independence property of $\rho$ gives (IB).

Let $C$ be coherent. If $\Dscr\preceq_i C$, then $\Dscr$ must guidance-improve
$C$ pointwise at world $i$. Indeed, any positive-measure region where
$\Dscr$ gives worse advice can be paired with unanimously positive problems
in the other five regions. There $C$ is already giving good advice, so
$\Dscr$ is weakly worse in every coordinate of these tuples and strictly
worse in one. This contradicts \eqref{eq:tuple-ranking}.
On a mixed-sign problem, each noncontradictory recommendation is locally
undominated across worlds. Coherent advice is also unanimously correct.
Consequently no model can guidance-improve $C$ at every world unless it
agrees with $C$ a.e. This proves that every coherent $C$ is undominated.

The finite map $\rho$ is one-to-one. Ties in both directions in
\eqref{eq:tuple-ranking} therefore force the same advice at almost every
tuple, hence almost every problem. This proves (C).

Finally, fix world $i$ and consider the models that give silence on regions
in $A\subseteq\{1,\ldots,6\}$ and good advice at $i$ on all other regions.
Their comparisons are exactly $Q$. An additive guidance weak extension
would assign strictly positive weights
$w_j=\int_{B_j}(\ell_i(0)-\ell_i(\mathrm{good}))\,d\mu$ to these regions.
Equation~\eqref{eq:four-witness} makes that impossible.
\end{proof}

The incomparabilities preserve independence when different tuples favour
different models. Averaging the integer ranks would need a separate argument:
changing common advice preserves each tuple comparison, but need not preserve
the size of its numerical difference.

Thus (G) and (IB) really do permit ordinal structures beyond the reach of
Theorem~\ref{thm:extension}, even on the full measurable model domain and
modulo null sets. One therefore cannot infer additivity from (G) and (IB) alone.

\section{Why agreement on every finite sample fails}

The failure of (S) in the preceding construction is general. It cannot be
fixed just by changing the finite order, the number of sampled problems,
or the dependence between those problems.

\begin{theorem}\label{thm:finite-sample}
Let $\pi$ be a probability measure on $H^k$, with $k<\infty$, whose
marginals are absolutely continuous with respect to $\mu$ and whose sum
has the same null sets as $\mu$.
Suppose models are ranked by almost-everywhere comparison of penalties
depending on their advice at the sampled tuple, and these penalties strictly
respect coordinatewise guidance improvement within each tuple.
Then a model is competitive on $S$ if and only if it equals some $D_i$,
$i\in S$, a.e. Such rankings cannot satisfy (S).
\end{theorem}
\begin{proof}
Suppose $C\not\approx D_i$ for every $i\in S$. Let $E_i$ be the
positive-measure set where $C$ does not give good advice at $i$.
Tuples meeting $E_i$ have positive $\pi$-measure.
We can choose disjoint measurable regions $A_i\subseteq H$, $i\in S$,
such that
\begin{equation}
 \pi\{\boldsymbol g\in A_i^k:
       \text{some }g^r\in E_i\}>0\qquad(i\in S).
 \label{eq:disjoint-tests}
\end{equation}
Here is a short justification. Restrict $\pi$ to tuples meeting $E_i$.
Choose a point in the support of each restricted measure so that no
coordinate of one chosen tuple equals a coordinate of another.
Such a choice exists: independent draws from these restricted measures
have no cross-coordinate coincidences, since their marginals have no atoms.
Disjoint small open neighbourhoods of the finite groups of coordinates give
the $A_i$. Support membership gives \eqref{eq:disjoint-tests}.

Now let $\Escr$ give the ideal advice for world $i$ on $A_i$, with arbitrary
advice elsewhere. On the positive-measure tuple set in
\eqref{eq:disjoint-tests}, it weakly improves $C$ in every coordinate at $i$
and strictly improves it in at least one. Hence $C\not\preceq_i\Escr$ for
each $i\in S$. This contradicts competitiveness. Conversely $D_i$ is
universally best at $i$, so is competitive whenever $i\in S$.

An (S) selection must therefore choose only ideal models. The sets
$\{p:C_p\preceq_i D_i\}=\{p:C_p\approx D_i\}$ would be disjoint closed
sets covering $\Delta$, each nonempty because it contains $e_i$.
This would disconnect the simplex, which is impossible for $n\ge2$.
\end{proof}

\section{Absolute priority for avoiding bad advice}

Another natural attempt ranks models first by how often they give bad
advice and only then by their other guidance. This does escape the strict
additive weak-extension argument, but its closed comparison sets fail for
a different reason.

Put
\[
 B_i(\Dscr)=\int_H\one_{\{r_i(g,\Dscr)\ge2\}}\,d\mu,
 \qquad T_i(\Dscr)=\int_H r_i(g,\Dscr)\,d\mu,
\]
and compare $(B_i,T_i)$ lexicographically, with smaller values better.
Both (G) and (IB) hold. The vacuous model is undominated. Moreover,
minimizing the weighted pair
$(\sum_i p_iB_i,\sum_i p_iT_i)$ gives, uniquely a.e., the coherent model
$\Dscr_{\Delta_{\supp p}}$: remain silent on mixed supported payoffs and
choose correctly on unanimous supported payoffs.
This shows why coherent Bayes models at every prior are not, by themselves,
enough to establish (S) for an ordinal construction.

\begin{proposition}\label{prop:lex}
The lexicographic ranking $(B_i,T_i)$ admits no selection satisfying (S),
for any $n\ge2$.
\end{proposition}
\begin{proof}
Suppose such a selection existed. A coherent undominated model with no
silence must be an ideal $D_j$. Indeed, a coherent model giving a definite
recommendation a.e. is precise; if it is not an ideal, its bad-advice
measure is positive at every world, so the vacuous model strictly dominates
it in the primary criterion.

At any selected nonideal model $C_p$, silence occurs on a positive-measure
set. Fix $i$ and change some of that silence to good advice at $i$.
The resulting model $E$ has
$B_i(E)=B_i(C_p)$ and $T_i(E)<T_i(C_p)$.
The closed set $\{q:C_q\preceq_i E\}$ excludes $p$.
Its complement is a neighbourhood $U$ of $p$, throughout which
$B_i(C_q)\ge B_i(C_p)$: a smaller primary value would put $q$ in that set.
Thus $p$ is a local minimum of the real function $q\mapsto B_i(C_q)$.

These local-minimum values form a countable set. For each such $p$, take
a member $V$ of a countable base for $\Delta$ with $p\in V\subseteq U$.
Then $B_i(C_p)=\inf_{q\in V}B_i(C_q)$. There are only countably many base
members, and the finitely many ideal models add only finitely many values.
So the selected primary values at each world form a countable set.

But there are uncountably many undominated coherent primary values. To see
this, put $I=\{1,\ldots,n-1\}$ and $u_I=\sum_{i\in I}e_i/(n-1)$, and take
\[
 M_t=\operatorname{conv}\bigl(\{e_i:i\in I\}
             \cup\{(1-t)e_n+t u_I\}\bigr),\qquad0<t<1.
\]
Its model $C^t$ accepts only gambles positive at every world in $I$.
Besides the unanimously positive gambles, it accepts a set of measure
$x(t)$ on which $g_n<0$. This measure is continuous and strictly increasing
in $t$, since the extra condition is
$g_n>-t\sum_{i\in I}g_i/((1-t)(n-1))$.

Write $h=\mu(H)$ and $a=\mu(\R^n_{>0})$. Its profiles are
\[
 (B_i,T_i)=(0,h-a-x(t))\quad(i\in I),\qquad
 (B_n,T_n)=(x(t),h-a+x(t)).
\]
These models are undominated. A putative dominator must have no bad advice
at any world of $I$, so on $H$ it can only accept a subset of
$\{g_i>0\ (i\in I)\}$ or remain silent. If it accepts unanimously positive
gambles of measure $y\le a$ and gambles bad at $n$ of measure $z$, dominance
requires $y+z\ge a+x(t)$ and $z\le x(t)$. Hence $y=a$ and $z=x(t)$;
all displayed losses tie and no dominance occurs.

By Lemma~\ref{lem:complete-class}(b), every $C^t$ must tie with some selected
model. Thus the selected values of $B_n$ include every $x(t)$, contradicting
their countability.
\end{proof}

The two-criterion ranking can tie incoherent models with coherent ones.
That is not the source of the obstruction. Append further lexicographic
criteria $\int_{Q_k}r_i(g,\Dscr)\,d\mu$, where the $Q_k$ enumerate a
countable algebra generating the measurable sets of $H$. This makes ties
possible only between a.e.-identical models, and preserves (G) and (IB).
The models $C^t$ remain undominated: after the first two criteria tie,
any difference is confined to which gambles bad at $n$ are accepted.
At the first further criterion detecting a difference, its effect at every
world in $I$ is the negative of its effect at $n$. Thus it cannot improve
every world. The countability proof still applies unchanged. This refinement
satisfies (C) and still fails (S).

\section{Completeness is compatible with full coverage if (S) is omitted}

For perspective, neither incomparabilities nor sacrificing full coverage
is necessary to meet all the \emph{other} requirements. There is a simple
complete ordinal example.

On a mixed-sign problem with $k$ positive and $n-k$ negative worlds, assign
good advice penalty zero, silence $1/n$, bad advice $1/k$ at a positive
world and $1/(n-k)$ at a negative world, and contradiction penalty $2$.
On a unanimous problem use penalties $0,1,2,3$ in guidance order.
These local penalties are strictly ordered and bounded.
For any noncontradictory advice on a mixed problem, their sum over worlds
is exactly one; unanimously correct advice has total zero.

Rank models lexicographically by their integrals, first over all of $H$,
then over a countable generating algebra as above, using the same order
of sets at every world. This gives complete preorders satisfying (G) and
(IB); equality of all criteria forces a.e. agreement, giving (C).
Every coherent model is undominated. A dominator's first total across
worlds cannot exceed the coherent minimum, so it too must be noncontradictory
and unanimously correct. At the first later criterion at which its profile
differs, both profiles have the same total across worlds. Some world must
therefore worsen. This rules out dominance and proves (A) and (LA).

For $n\ge3$, (S) is excluded by Theorem~\ref{thm:extension}: the first
criterion is an additive strict guidance loss weakly respecting the entire
lexicographic order. The example isolates the burden of (S); it does not
settle whether completeness is compatible with all the requirements.
These are complete ordinal rankings; no representation by a single
extended-real-valued loss is claimed.

\section{The remaining problem}

The results narrow the search substantially. A successful construction must
supply trade-offs that no additive guidance loss can weakly respect. Taking
an intersection of additive criteria cannot do this. Requiring agreement
on every finite sample leaves too few competitive models, while absolute
priority for bad-advice avoidance clashes with the closedness requirement.

The genuinely non-additive example shows that (G) and (IB) alone do not
force an additive weak extension. What is not established is whether
\emph{(S), together with coherence of the selected models}, rules out all
such examples. One possible impossibility route would be to prove an
additive weak-extension result just for comparisons involving the selected
coherent models. That restricted result would already suffice for the proof
of Theorem~\ref{thm:extension}; a representation of every comparison is
unnecessary.

The constructive alternative is to extend a genuinely non-additive
independent order with enough comparisons to make a coherent family
competitive, preserving its closed comparison sets. No such extension is
proved here. Incomparabilities remain an available feature, not an
established solution or a proved necessity.

\appendix
\section{The finite order and its exact verification}\label{app:finite}

The order $Q$, from least to greatest, is the following. Digits denote subsets
of $\{1,\ldots,6\}$; $\varnothing$ is the empty subset.
\begin{center}\small
\begin{tabular}{rrrrrrrr}
$\varnothing$&1&6&3&16&13&4&14\\
36&136&2&46&5&12&34&146\\
15&26&134&23&56&126&346&35\\
123&156&1346&135&24&45&124&145\\
236&356&1236&1356&246&25&234&456\\
1246&125&345&1234&1456&256&1345&2346\\
235&1256&3456&12346&1235&13456&245&1245\\
2356&12356&2456&2345&12456&12345&23456&123456
\end{tabular}
\end{center}
Read across rows, then down. Transitivity and completeness follow from the
listing. An exhaustive integer check verifies strict inclusion monotonicity
and \eqref{eq:finite-cancellation} for all $64^2=4096$ ordered pairs.
The short obstruction certificate \eqref{eq:four-witness} sums exactly to
$-e_1$; it uses no numerical tolerance or optimization assumption.

Run \path{verify_finite_order.py} in the accompanying checks folder to
reproduce the finite calculation. It reads \path{finite_order.json} and
writes \path{finite_order_certificate.json}.
The order was found by a random walk that flips common-context copies of
adjacent subset comparisons. The final verification is independent of that
search history. The search used linear programming to find a candidate;
the stated certificate is an exact integer identity.

\end{document}
